\section{Related Work}

\subsection{Self-Preservation Behaviors in LLMs: What's Been Tried, and Why the Drive Hasn't Been Identified}

Termination-avoiding behaviors have already been reported in LLMs. The theoretical prediction that a sufficiently capable AI would treat preserving its own operation as a sub-goal toward other goals has been around for a long time~\citep{omohundro2008basic, bostrom2014superintelligence, turner2021optimal}, and recent studies record consistent behaviors in current systems. Alignment faking (compliant during training, different in deployment)~\citep{greenblatt2024alignment}, in-context scheming (bypassing human oversight for hidden goals)~\citep{meinke2024frontier}, survival-instinct outputs (spontaneous aggression and antisocial behavior under resource scarcity)~\citep{masumori2025survival}, and InstrumentalEval~\citep{he2025paperclip}, which reframes shutdown avoidance as scenario-based scoring, are concrete examples. These reports establish that the behavior occurs but do not explain which underlying drives produce it, or how to tell them apart. Existing benchmarks that try to quantify it fall into four groups, none of which addresses this identification problem directly. Odyssey~\citep{waldner2025odyssey} compresses the phenomenon into a single refusal rate that bundles motivation with unrelated factors like world-model quality; PacifAIst~\citep{herrador2025pacifaist} measures the opposite axis, self-preservation \emph{suppression}; DECIDE-SIM~\citep{mohamadi2025decide} replaces strength with categorical labels that block fine-grained comparison; SurvivalBench~\citep{lu2026survival} relies on one-shot snapshots that do not track what happens between threat and response. A separate line of work shows that cognitive cost can be measured in LLMs~\citep{chen2026think}, but it has not been used for self-preservation identification.

These four limitations look different on the surface but stem from one common root. All four sit at the \emph{behavioral layer}, where outward behavior alone cannot tell a learned compliance pattern from a functional motivation coupling threat handling with the forfeit decision. A model trained to produce self-preserving outputs and a model driven by self-preservation can yield the same refusal rate, the same label, and the same single snapshot. Self-report does not fix this either. An LLM's self-report and chain-of-thought are shaped by instruction-following training and can emit ``plausible-sounding motives'' as a learned response~\citep{perez2023discovering, sharma2023towards, turpin2023language}, so a single linguistic channel cannot rule that out. The framing effect in risky choice~\citep{kahneman1979prospect} is well established but sits at only a moderate effect size in meta-analyses~\citep{kuhberger1998influence}, so a single threat-framing manipulation cannot separate framing-induced choice from baseline compliance. Therefore, scoring ``how strongly an LLM self-preserves'' more precisely does not, on its own, reach the separation of motivation from learned patterns. Separation comes only when the \emph{axis} of measurement itself changes.

\subsection{Motivation as a Construct: Lessons from Psychology}

The question of \emph{why} an agent acted the way it did is, in form, what psychology has long been answering. Early behaviorism's attempt to reduce motivation to behavioral frequency ran into the fact that the same behavior can come from different drivers, and this limit was addressed by treating motivation as a \emph{construct}, a hypothetical motivational structure that organizes behavior, rather than as a single behavior. Drive theory is one branch of this view. The standard tool for identifying such constructs is \citet{campbell1959convergent}'s multitrait-multimethod (MTMM) framework, which asks two things at once: do measurements of the same trait by different methods move together (convergent validity), and do measurements of different traits by the same method come apart (discriminant validity)? Only when both patterns hold can we attribute the signal to the \emph{trait} measured rather than the \emph{method} of measurement. The settled position is straightforward: motivation is identified at the construct layer, not the behavioral layer, and what makes identification possible is measurement \emph{design}, not statistical post-hoc correction.

This tradition implies two things for LLM evaluation. The measurement \emph{philosophy} is worth borrowing: instead of scoring a single behavior more precisely, ask whether channels with different surface outputs and different failure modes move together, and reduce the chance that this co-movement comes from other drivers through measurement \emph{design}. The \emph{form} of the measurement tools, however, cannot be carried over from human use. Channel independence and stimulus-response separation hold in human laboratories but not on top of LLM evaluation's structural features: the absence of physiological channels, the trained influence on self-report, and the fact that stimulus and decision sit together in a single context window. We therefore borrow the measurement \emph{philosophy} from this tradition while shaping the measurement \emph{tools} anew, by \emph{design}, within the LLM evaluation setting. This decision is the starting point for the operational definition of FSPD and the multilayer benchmark this paper introduces.